% Fonte LaTeX do artigo. Substitua o preambulo pelo template SBC quando necessario.
\documentclass[10pt,twocolumn]{article}
\usepackage[T1]{fontenc}
\usepackage[utf8]{inputenc}
\usepackage[brazil]{babel}
\usepackage{graphicx}
\usepackage{booktabs}
\usepackage{array}
\usepackage{url}
\usepackage{hyperref}
\usepackage{amsmath}
\usepackage[margin=2cm]{geometry}
\title{Deteccao de Fraudes em Cartoes com Lakehouse Local, Apache Spark e Avaliacao Temporal Reprodutivel}
\author{Rafael Luciano}
\date{Junho de 2026}

\begin{document}
\maketitle

\begin{abstract}
A deteccao de fraude em pagamentos combina uma classe positiva rara, custos assimetricos de erro e necessidade de resposta rapida. Este trabalho apresenta e avalia uma arquitetura lakehouse local para o problema, integrando Kafka, Apache Spark Structured Streaming, HDFS, Apache Iceberg, Hive Metastore, Airflow e Spark MLlib. O experimento usa as 284.807 transacoes do conjunto ULB e separa cronologicamente treino, validacao e teste em 60/20/20. Dois ensembles Random Forest--Gradient-Boosted Trees ponderados por classe sao comparados em cinco sementes; a selecao usa somente PR-AUC de validacao e o teste permanece isolado. O candidato selecionado alcancou, no bloco de teste, PR-AUC de $0{,}749825 \pm 0{,}016599$, ROC-AUC de $0{,}978700 \pm 0{,}004384$, precision de $0{,}861579 \pm 0{,}025262$ e recall de $0{,}735135 \pm 0{,}015408$. A materializacao operacional produziu 284.807 scores e 562 alertas. Alem de resultados preditivos, o estudo mostra como uma camada de tabelas transacionais e um catalogo compartilhado sustentam rastreabilidade, reprocessamento e analise de qualidade.
\end{abstract}

\textbf{Palavras-chave:} deteccao de fraude; lakehouse; Apache Spark; streaming; Apache Iceberg; avaliacao temporal; dados desbalanceados.

\section{Introduction}

Pagamentos digitais tornaram a analise antifraude uma tarefa simultaneamente estatistica e operacional. A decisao precisa ocorrer no intervalo entre a chegada do evento e sua autorizacao, mas o resultado verdadeiro da transacao normalmente so e conhecido depois. Alem disso, uma pequena fracao dos registros representa fraude. Uma acuracia elevada, portanto, pode ser obtida por um classificador inutil que chama todos os eventos de legitimos. A consequencia pratica e dupla: falsos negativos expoem a instituicao a perdas; falsos positivos interrompem clientes legitimos e aumentam o custo de revisao humana.

A literatura de fraude enfatiza que a modelagem deve considerar desbalanceamento, mudanca temporal e custos de decisao, em vez de tratar a classificacao binaria como um problema i.i.d. comum \cite{ngai2011,dalpozzolo2015,dalpozzolo2018}. Em particular, um teste cronologicamente posterior e mais proximo do uso operacional do que uma particao aleatoria: ele reduz a possibilidade de registros de um mesmo periodo aparecerem nos dois lados da avaliacao e oferece uma visao mais honesta da generalizacao ao futuro. Essa preocupacao e especialmente relevante quando as variaveis descrevem padroes transacionais, pois comportamento fraudulento e comportamento legitimo podem evoluir.

Ao lado do modelo, a infraestrutura precisa preservar o caminho do dado. Uma solucao que recebe eventos em streaming, treina em lote e depois produz indicadores analiticos tende a acumular sistemas separados, copias de dados e regras dificeis de auditar. O conceito de lakehouse procura combinar armazenamento aberto de um data lake com controle de metadados, confiabilidade de tabelas e desempenho analitico tradicionalmente associados ao data warehouse \cite{armbrust2021}. Essa convergencia e adequada a fraude: os mesmos dados podem sustentar ingestao, preparacao, treinamento, inferencia, alertas e consultas retrospectivas.

Este artigo apresenta uma implementacao local, deliberadamente reprodutivel, de um lakehouse para deteccao de fraude. A contribuicao principal nao e propor um novo algoritmo, mas estabelecer uma pipeline experimental completa e auditavel com ferramentas abertas. O trabalho: (i) organiza eventos ULB em camadas Bronze, Silver e Gold sobre Iceberg/HDFS; (ii) realiza inferencia apos Silver e antes de Gold, preservando a ordem operacional; (iii) define um protocolo temporal de treino, validacao, teste e cinco repeticoes; e (iv) reporta metricas de classificacao, matrizes de confusao, scores, alertas, qualidade e tempos de execucao. Trabalhos recentes sobre lakehouse para risco de pagamentos reforcam a utilidade de combinar processamento de eventos e analise em lote, embora frequentemente descrevam arquiteturas em nivel conceitual \cite{battapothu2025,carrington2025,adlermann2024}. Aqui, cada resultado apresentado foi materializado no ambiente executado.

\section{Description} \label{sec:firstpage}

\subsection{Deteccao de fraude e avaliacao desbalanceada}

Em deteccao de fraude, a classe positiva e rara e os erros nao tem o mesmo significado. Seja $TP$ o numero de fraudes corretamente sinalizadas, $FP$ os alertas incorretos, $TN$ as transacoes legitimas aceitas e $FN$ as fraudes nao sinalizadas. A precision e $TP/(TP+FP)$, o recall e $TP/(TP+FN)$, enquanto o F1 harmoniza os dois criterios. Essas medidas devem ser lidas em conjunto: aumentar a sensibilidade reduzindo o limiar pode elevar fortemente os falsos positivos; tornar o limiar muito alto pode reduzir a fila de revisao, mas deixar fraudes passar.

A curva ROC e seu AUC continuam uteis para resumir ordenacao de scores, mas podem parecer excessivamente favoraveis quando negativos dominam a populacao. A curva precision--recall (PR) concentra a analise na classe de interesse e e recomendada para conjuntos desbalanceados \cite{davis2006,saito2015}. Por isso, este estudo declara PR-AUC como metrica primaria para escolher hiperparametros. ROC-AUC, precision, recall, F1 e a matriz de confusao sao reportados como medidas complementares. A acuracia nao e utilizada para selecao.

Outra decisao fundamental e a separacao dos dados. Validacao cruzada estratificada e comum em problemas estaticos, mas pode subestimar o erro quando ha dependencia temporal. Estudos de fraude recomendam distinguir o periodo de treinamento do periodo futuro de avaliacao e tornar explicita a estrategia de atualizacao \cite{dalpozzolo2018}. O protocolo deste trabalho usa tres blocos cronologicos e nao permite que o teste influencie a escolha de algoritmo, limiar ou pesos.

\subsection{Lakehouse e tabelas transacionais abertas}

O lakehouse emprega armazenamento de arquivos aberto, mas adiciona um catalogo e operacoes de tabela para oferecer consistencia a workloads de engenharia e analise. Iceberg materializa esse papel no experimento: cada camada e uma tabela catalogada pelo Hive Metastore e mantida no HDFS. O catalogo fornece um nome estavel para leitura por Spark, enquanto snapshots e metadados permitem evolucao e inspeccao das tabelas. Dessa forma, a camada Bronze nao e apenas um diretorio bruto, e Gold nao e apenas uma exportacao CSV sem linhagem.

As camadas seguem a convencao medallion. Bronze retém evento, topico, particao, offset, instante de ingestao e identificador tecnico. Silver aplica tipos, remove registros invalidos, cria a label e disponibiliza as features do modelo. Gold recebe scores, alertas, KPIs, metricas de qualidade e resultados experimentais. A separacao torna visivel onde um resultado foi produzido e permite reprocessar Silver ou Gold sem republicar o arquivo de origem. A visao e compativel com arquiteturas multi-tier discutidas para dominios financeiros regulados \cite{bfsilakehouse2023,allam2024}.

\subsection{Processamento distribuido, eventos e orquestracao}

Apache Spark e um motor unificado para processamento em memoria, SQL, streaming estruturado e aprendizado de maquina \cite{zaharia2016}. No projeto, o mesmo ecossistema executa a limpeza batch, o treinamento MLlib, a leitura Structured Streaming de Kafka, a gravacao Iceberg e as agregacoes Gold. Isso reduz a troca de formatos entre etapas e permite reutilizar o catalogo e as mesmas definicoes de schema.

Kafka e o barramento de eventos. O produtor publica uma transacao ULB por mensagem JSON em \texttt{transactions.raw}; o streaming Bronze a le a partir do inicio e persiste offsets no checkpoint HDFS. Kafka foi concebido como sistema de mensagens distribuido para processamento de logs, com retencao e alto throughput \cite{kreps2011}. Structured Streaming oferece o modelo de micro-batches e checkpoint para que a transformacao seja expressa como consulta incremental, aproximando processamento de eventos e processamento historico \cite{akidau2015}.

Airflow tem papel de orquestrador: suas DAGs disparam os jobs de inicializacao, ingestao, Silver, treinamento, agregacao e qualidade. Ele nao controla um loop continuo de consumo Kafka; os streams pertencem ao Spark. Essa separacao evita usar um agendador batch como mecanismo de streaming e preserva a responsabilidade de cada componente. Artigos aplicados recentes tambem relacionam orquestracao e camadas analiticas a rastreabilidade em ambientes financeiros \cite{bfsilakehouse2023}.

\section{Methodology}

\subsection{Visao geral da arquitetura}

A arquitetura foi implementada inteiramente em contêineres Docker Compose. Os servicos sao Kafka e ZooKeeper, HDFS NameNode/DataNode, Hive Metastore com PostgreSQL, Spark Master e Worker, Airflow com PostgreSQL e Kafka UI. O armazenamento principal e HDFS; Iceberg usa Hive Catalog no endereco do metastore. As versoes sao fixadas no repositório: Spark 3.5.1, Iceberg 1.5.2, Kafka Confluent 7.6.1, Hadoop 3.2.1, Hive Metastore 3.1.3 e Airflow 2.9.2.

\begin{center}
\fbox{\parbox{0.92\columnwidth}{\centering \textbf{[Figura 1 -- inserir manualmente]}\\ Diagrama da arquitetura: CSV ULB $\rightarrow$ produtor $\rightarrow$ Kafka $\rightarrow$ Bronze/Iceberg-HDFS $\rightarrow$ Silver $\rightarrow$ treinamento MLlib $\rightarrow$ scoring streaming ou batch $\rightarrow$ Gold, Power BI e Airflow.}}
\end{center}

O desenho atende a dois caminhos. O caminho online recebe e persiste eventos, enquanto o caminho analitico produz modelos, analisa resultados e prepara indicadores. O score e calculado depois da Silver, pois somente ali as features tem schema e tipos controlados; ele ocorre antes da Gold, pois Gold armazena o produto da decisao, e nao a etapa de decisao. Essa ordem tambem permite que alertas e scores sejam publicados no Kafka sem aguardar agregacoes diarias.

\subsection{Ingestao e camada Bronze}

O arquivo \texttt{creditcard.csv} e validado antes da publicacao. O produtor verifica a presenca de \texttt{Time}, \texttt{V1} a \texttt{V28}, \texttt{Amount} e \texttt{Class}; cada linha recebe um identificador deterministico e e enviada a Kafka. Na execucao, foram usados lotes de 5.000 mensagens. O produtor publicou 284.807 eventos em 164,6 s, distribuidos em tres particoes do topico. O checkpoint Bronze confirmou os offsets finais 95.515, 95.267 e 94.025, cuja soma e o numero total de eventos.

O job \texttt{kafka\_to\_bronze.py} le JSON, preserva o offset Kafka e grava em \texttt{local.bronze.transactions\_raw}. A combinacao de \texttt{transaction\_id} e offset cria \texttt{raw\_event\_id}, permitindo distinguir replay tecnico de identidade de negocio. O checkpoint fica no HDFS; assim, uma reinicializacao do stream continua a partir do ponto processado. Essa etapa foi mantida mesmo para o conjunto historico porque o objetivo experimental inclui validar a rota de ingestao que seria usada para dados em producao.

\subsection{Silver: qualidade, schema e features}

Silver e gerada pelo job \texttt{bronze\_to\_silver.py}. Todas as variaveis do dataset sao convertidas explicitamente para \texttt{double}; a classe e disponibilizada como \texttt{label}; valores nulos das features sao preenchidos com zero e labels fora de $\{0,1\}$ sao excluidos. A tabela \texttt{local.silver.transactions\_clean} preserva metadados tecnicos e \texttt{local.silver.transactions\_features} inclui um vetor de features para consumo analitico.

O desbalanceamento e tratado por pesos de classe no treinamento. Para cada classe $c$, o peso empregado e $w_c=N/(2N_c)$, em que $N$ e o numero de registros de treinamento e $N_c$ o numero de elementos da classe. Os pesos sao calculados exclusivamente no bloco de treino em cada etapa de ajuste, nao no conjunto completo. Essa estrategia preserva o conjunto original e evita introduzir sintetizacao externa, alinhando-se a uma comparacao controlada com algoritmos Spark MLlib.

\subsection{Protocolo experimental e prevencao de vazamento}

O protocolo foi registrado em \texttt{configs/app/experiment.yaml}. A variavel \texttt{Time} determina dois pontos de corte aproximados por quantil: 120.390 s e 145.229 s. O primeiro bloco contem 170.229 transacoes (60\%), o segundo 56.732 (20\%) e o bloco de teste 56.765 (20\%). Os tamanhos diferem levemente das proporcoes nominais porque existem eventos com o mesmo instante; a regra de fronteira preserva o ordenamento. Os blocos contem, respectivamente, 323, 95 e 74 fraudes.

O experimento adota cinco sementes: 42, 52, 62, 72 e 82. Nao existe um unico numero universal de repeticoes que seja ``gold standard'' para fraude; a literatura salienta que a avaliacao deve refletir a dependencia temporal e a variabilidade do metodo \cite{dalpozzolo2018,hand2009}. Cinco repeticoes sao usadas aqui como compromisso usual de estabilidade para algoritmos estocasticos em ambiente local: permitem reportar media e desvio-padrao, sem confundir o teste com uma busca exaustiva de configuracoes. O conjunto de teste permanece fechado durante as dez execucoes de validacao.

Dois candidatos foram definidos antes da execucao. O candidato \textit{compact} usa Random Forest com 80 arvores e profundidade maxima 8, mais GBT com 50 iteracoes e profundidade 3. O candidato \textit{extended} usa 120 arvores/profundidade 10 e GBT de 60 iteracoes/profundidade 4. A grade e deliberadamente pequena: arvores profundas e muitas iteracoes aumentariam custo e risco de ajuste excessivo para uma demonstracao arquitetural, enquanto duas capacidades permitem observar se maior complexidade melhora a ordenacao PR no periodo de validacao. Cada candidato e treinado cinco vezes sobre o mesmo bloco de treino.

Para cada ajuste, o score do ensemble e a media simples das probabilidades de fraude de Random Forest e GBT:
\begin{equation}
s(x) = \frac{p_{RF}(y=1|x) + p_{GBT}(y=1|x)}{2}.
\end{equation}
O limiar e buscado na grade $0{,}05, 0{,}10, \ldots, 0{,}95$. A politica seleciona o maior recall com precision minima de 0,70; se nenhum limiar satisfizer o piso, usa o maior F1. A decisao de configuracao usa a media de PR-AUC na validacao, com F1 como desempate. Apenas apos a escolha sao carregados os cinco modelos vencedores e calculadas as metricas no teste. Por fim, um modelo de implantacao e ajustado sobre treino mais validacao; seu limiar e a mediana dos cinco limiares de validacao. Esse ultimo modelo materializa Gold, mas suas metricas retrospectivas nao substituem o teste isolado.

\subsection{Treinamento, registro e inferencia}

O treinamento usa exclusivamente \texttt{RandomForestClassifier} e \texttt{GBTClassifier} do Spark MLlib, ambos com \texttt{weightCol}. Cada modelo de validacao e salvo no HDFS em um caminho que contem identificador do experimento, candidato e semente. As tabelas \texttt{local.gold.experiment\_metrics} e \texttt{local.gold.experiment\_summary} registram seed, split, limiar, matriz de confusao, PR-AUC, ROC-AUC, media e desvio-padrao. O arquivo local \texttt{models/latest\_metadata.json} registra versao, colunas, pesos, pontos de corte e caminhos dos modelos.

Para demonstrar materializacao reprodutivel, o job \texttt{score\_silver\_batch.py} aplica o modelo de implantacao a toda a Silver. O job de streaming existente aplica a mesma logica em micro-batches para novos registros. Em ambos os casos, \texttt{fraud\_score}, probabilidades individuais, limiar, versao do modelo e instante de score sao escritos em \texttt{local.gold.fraud\_scores}; previsoes positivas tambem entram em \texttt{local.gold.fraud\_alerts}. A rotina Gold gera KPIs horarios e diarios e exporta CSV para Power BI. A rotina de qualidade mede contagem, nulidade, distribuicoes e outros indicadores, gravando 151 metricas na tabela Gold.

\subsection{Ambiente e verificacao}

O experimento foi executado em 21 de junho de 2026 no ambiente local Docker do projeto, com um Spark Worker configurado com 2 cores e 3 GB de memoria. A execucao completa incluiu reinicializacao dos volumes para evitar mistura com a amostra exploratoria anterior, inicializacao Iceberg, criacao dos topicos, ingestao, Bronze, Silver, experimento, scoring batch, Gold e qualidade. Os testes unitarios do repositorio foram executados no container Spark e obtiveram 12 testes aprovados.

Os tempos observados foram 164,6 s para publicacao Kafka, 66,7 s para Bronze--Silver, 1.224,8 s para o experimento completo de modelos, 74,0 s para scoring batch, 44,7 s para agregacoes Gold e 68,2 s para qualidade. Esses tempos nao sao benchmark de escala, pois incluem um unico worker e downloads iniciais de dependencias no primeiro job. Eles documentam, contudo, que a pipeline foi executada de ponta a ponta e oferecem uma base para comparacao de futuras configuracoes.

\section{Results}

\subsection{Dataset ULB e integridade da execucao}

O conjunto Credit Card Fraud Detection foi disponibilizado pelo Machine Learning Group da Universite Libre de Bruxelles em parceria com Worldline \cite{ulb2013}. Ele contem transacoes de cartoes europeus realizadas durante dois dias de setembro de 2013. Das 284.807 instancias, 492 sao fraudulentas, o que corresponde a aproximadamente 0,1727\%. Por razoes de confidencialidade, as componentes \texttt{V1}--\texttt{V28} resultam de transformacao PCA; \texttt{Time} representa segundos desde a primeira transacao e \texttt{Amount} e o valor monetario. A label \texttt{Class} e usada apenas para treinamento e avaliacao offline; em uma implantacao real, o rotulo chegaria com atraso.

A Tabela \ref{tab:layers} confirma a integridade basica do fluxo. Bronze, Silver e Gold de scores preservaram as 284.807 transacoes. Isso mostra que o produtor, as tres particoes Kafka, o stream Iceberg e o job Silver nao perderam ou duplicaram registros no experimento reinicializado. A tabela de alertas contem apenas a parcela classificada como suspeita; qualidade registrou 151 medidas auditaveis.

\begin{table}[ht]
\centering
\caption{Contagens materializadas por camada.}
\label{tab:layers}
\begin{tabular}{lr}
\toprule
Artefato & Registros \\
\midrule
Bronze \texttt{transactions\_raw} & 284.807 \\
Silver \texttt{transactions\_clean} & 284.807 \\
Gold \texttt{fraud\_scores} & 284.807 \\
Gold \texttt{fraud\_alerts} & 562 \\
Metricas de qualidade & 151 \\
\bottomrule
\end{tabular}
\end{table}

\subsection{Selecao por validacao}

A Tabela \ref{tab:validation} mostra as medias das cinco sementes. O candidato \textit{compact} foi escolhido porque sua PR-AUC media na validacao foi 0,656087, ligeiramente superior a 0,649046 do \textit{extended}. A diferenca e pequena e os desvios-padrao se sobrepoem; portanto, o resultado nao deve ser lido como evidencia de superioridade ampla da configuracao menor. Ele mostra, no entanto, que aumentar arvores, profundidade e iteracoes nao trouxe ganho consistente no criterio previamente declarado. O \textit{extended} obteve precision e F1 medias superiores nesse bloco, mas a regra de selecao manteve PR-AUC como objetivo principal para evitar trocar o criterio apos ver os resultados.

\begin{table*}[t]
\centering
\caption{Validacao temporal: media $\pm$ desvio-padrao em cinco sementes. O vencedor foi definido por PR-AUC.}
\label{tab:validation}
\begin{tabular}{lrrrrrr}
\toprule
Candidato & PR-AUC & ROC-AUC & Precision & Recall & F1 & Limiar \\
\midrule
Compact & $0{,}656087\pm0{,}030716$ & $0{,}972292\pm0{,}002404$ & $0{,}797130\pm0{,}049850$ & $0{,}736842\pm0{,}021487$ & $0{,}765472\pm0{,}032498$ & $0{,}70\pm0{,}00$ \\
Extended & $0{,}649046\pm0{,}045177$ & $0{,}959786\pm0{,}008407$ & $0{,}860991\pm0{,}063378$ & $0{,}715789\pm0{,}031383$ & $0{,}780548\pm0{,}032860$ & $0{,}58\pm0{,}027386$ \\
\bottomrule
\end{tabular}
\end{table*}

O limiar 0,70 foi selecionado em todas as cinco execucoes do candidato compacto. Isso e coerente com a politica de precision minima 0,70: no periodo de validacao, o modelo pode elevar recall ate o ponto em que a fila de alertas ainda preserva qualidade. O limiar do modelo de implantacao foi a mediana desses valores, tambem 0,70. A estabilidade do limiar nao elimina a necessidade de recalibracao futura, pois o dataset e historico e nao representa deriva de conceito em longo prazo.

\begin{center}
\fbox{\parbox{0.92\columnwidth}{\centering \textbf{[Figura 2 -- inserir manualmente]}\\ Curvas precision--recall das cinco sementes na validacao e no teste, com destaque para o candidato compacto selecionado.}}
\end{center}

\subsection{Teste cronologico isolado}

O bloco de teste contem 56.765 transacoes, das quais 74 sao fraudes. A Tabela \ref{tab:test} apresenta cada repeticao do candidato selecionado. A media de PR-AUC foi 0,749825 com desvio-padrao 0,016599; ROC-AUC foi 0,978700 com desvio-padrao 0,004384. Precision media de 0,861579 significa que a maior parte dos alertas do teste correspondeu a fraude, enquanto recall medio de 0,735135 indica que cerca de tres quartos das 74 fraudes foram sinalizadas no limiar escolhido. O F1 medio foi 0,793038.

\begin{table*}[t]
\centering
\caption{Resultados por semente no bloco temporal de teste.}
\label{tab:test}
\begin{tabular}{rrrrrrrrrrrr}
\toprule
Run & Seed & Precision & Recall & F1 & PR-AUC & ROC-AUC & TP & FP & TN & FN \\
\midrule
1 & 42 & 0,861538 & 0,756757 & 0,805755 & 0,749979 & 0,974861 & 56 & 9 & 56.682 & 18 \\
2 & 52 & 0,843750 & 0,729730 & 0,782609 & 0,756951 & 0,979404 & 54 & 10 & 56.681 & 20 \\
3 & 62 & 0,898305 & 0,716216 & 0,796992 & 0,721545 & 0,973738 & 53 & 6 & 56.685 & 21 \\
4 & 72 & 0,870968 & 0,729730 & 0,794118 & 0,764267 & 0,981293 & 54 & 8 & 56.683 & 20 \\
5 & 82 & 0,833333 & 0,743243 & 0,785714 & 0,756382 & 0,984204 & 55 & 11 & 56.680 & 19 \\
\midrule
Media & -- & 0,861579 & 0,735135 & 0,793038 & 0,749825 & 0,978700 & -- & -- & -- & -- \\
DP & -- & 0,025262 & 0,015408 & 0,009232 & 0,016599 & 0,004384 & -- & -- & -- & -- \\
\bottomrule
\end{tabular}
\end{table*}

Os resultados mostram baixa variacao em recall e F1, e variacao moderada em PR-AUC. Essa variacao e esperada porque ha apenas 74 positivos no teste e porque Random Forest e GBT possuem componentes estocasticos. O ponto mais baixo de PR-AUC ocorreu na semente 62, embora sua precision tenha sido a maior; novamente, uma unica metrica nao descreve todos os trade-offs. A avaliacao por repeticao e mais informativa que reportar somente o melhor run, pratica que poderia superestimar a capacidade do metodo.

O teste tambem evidencia um ponto de interpretacao: PR-AUC no teste foi maior que na validacao. Isso nao deve ser tratado como melhoria causada pelo teste, mas como diferenca de dificuldade entre janelas temporais e tamanho pequeno da classe positiva. Como o teste nao participou da escolha de configuracao, ele continua valido para a estimativa deste protocolo. A diferenca reforca a importancia de executar reavaliacoes temporais sucessivas em cenarios reais.

\subsection{Scores e alertas do modelo de implantacao}

Depois de concluida a avaliacao, o modelo compacto de implantacao foi treinado com treino e validacao e aplicado retrospectivamente a toda a Silver. Esta etapa nao e uma nova medida de generalizacao, porque inclui registros vistos no treino; seu objetivo e demonstrar o produto operacional da pipeline. Foram persistidos 284.807 scores e 562 alertas a limiar 0,70. Dos alertas, 435 eram fraudes reais e 127 eram falsos positivos; 57 fraudes nao foram alertadas. A precision retrospectiva foi $435/562=0,7731$ e o recall retrospectivo foi $435/492=0,8841$.

O recall operacional maior que o recall de teste e esperado porque o modelo de implantacao usa 80\% dos dados para treinamento e e aplicado tambem sobre esse periodo. Por essa razao, a Tabela \ref{tab:test}, e nao estes numeros retrospectivos, e a referencia cientifica para desempenho preditivo. Os scores Gold, contudo, sao valiosos para auditoria, priorizacao e BI: cada linha preserva probabilidades de RF e GBT, score final, predicao, limiar, versao do modelo e timestamp de scoring.

\begin{table}[ht]
\centering
\caption{Materializacao operacional sobre toda a Silver. Nao e avaliacao independente.}
\label{tab:operational}
\begin{tabular}{lr}
\toprule
Medida & Valor \\
\midrule
Scores produzidos & 284.807 \\
Fraudes rotuladas & 492 \\
Alertas & 562 \\
Verdadeiros positivos & 435 \\
Falsos positivos & 127 \\
Falsos negativos & 57 \\
Precision retrospectiva & 0,7731 \\
Recall retrospectivo & 0,8841 \\
\bottomrule
\end{tabular}
\end{table}

\begin{center}
\fbox{\parbox{0.92\columnwidth}{\centering \textbf{[Figura 3 -- inserir manualmente]}\\ Distribuicao de \texttt{fraud\_score} e linha vertical no limiar 0,70; opcionalmente sobrepor fraudes reais e alertas.}}
\end{center}

\subsection{Desempenho da pipeline e qualidade}

A execucao confirma que a arquitetura suporta o fluxo completo no ambiente local. A publicacao Kafka levou 164,6 s e o experimento, que inclui dez ajustes de validacao mais o modelo de implantacao, levou 1.224,8 s (aproximadamente 20,4 min). O scoring batch final levou 74,0 s. Esses numeros sao observacoes do ambiente de teste, e nao afirmacoes de throughput universal: ha apenas um worker, dois cores e 3 GB de memoria. Ainda assim, o fato de todo o conjunto atravessar Kafka, Iceberg e Spark sem perda fornece evidencia de integracao funcional.

As 151 metricas de qualidade incluem contagens, completude e distribuicoes por tabela. A igualdade das contagens entre Bronze, Silver e Gold scores e uma verificacao simples e forte para esta execucao. Em producao, esse conjunto deveria ser complementado por alertas de deriva de distribuicao, atraso de label, mudanca de precision estimada e SLA por micro-batch. A estrutura Iceberg permite que essas verificacoes sejam escritas como tabelas, em vez de logs efemeros.

\section{Conclusions}

Este trabalho desenvolveu e executou uma arquitetura lakehouse local para deteccao de fraude em cartoes. A pipeline integrou Kafka para eventos, Spark Structured Streaming para Bronze e inferencia, HDFS e Iceberg para persistencia de camadas, Hive Metastore para catalogo, Airflow para orquestracao e Spark MLlib para treinamento. O experimento ULB percorreu todas as 284.807 transacoes, gerou 284.807 scores, 562 alertas, KPIs, exportacoes CSV e 151 metricas de qualidade.

Do ponto de vista metodologico, a principal decisao foi substituir uma divisao aleatoria simples por protocolo temporal 60/20/20, com validacao separada, teste isolado e cinco sementes. O candidato compacto foi selecionado por PR-AUC de validacao e obteve PR-AUC medio de 0,749825 e precision media de 0,861579 no teste. Os valores mostram que o modelo identifica uma parcela relevante das fraudes mantendo uma fila de alertas relativamente precisa, mas tambem deixam claro que ainda ha falsos negativos e que a escolha do limiar depende do custo real de revisao e perda.

Ha limitacoes importantes. O ULB e historico, tem features anonimizadas por PCA e representa somente dois dias; logo, nao demonstra deriva de conceito, disponibilidade tardia de labels ou comportamento de clientes em longo prazo. O experimento utiliza um unico worker local e nao mede escalabilidade horizontal, latencia ponta a ponta por evento ou disponibilidade sob falha. A grade de hiperparametros foi pequena, os modelos nao incluem variaveis de grafo ou contexto de entidade e o ensemble por media nao substitui uma analise de custo monetario. Tambem nao foi realizado teste estatistico de diferenca entre configuracoes, pois os cinco runs compartilham os mesmos blocos temporais.

Como trabalhos futuros, propoe-se avaliar janelas temporais deslizantes e prequential, incorporar custos por transacao e limites de capacidade de revisao, comparar diferentes politicas de recalibracao, medir throughput com multiplos workers e adicionar monitoramento de deriva. A extensao com atributos de conta, comerciante e dispositivo permitiria investigar modelos relacionais, desde que preservada a separacao entre avaliacao e dados futuros. O estudo pode ser reproduzido pelo codigo no GitHub em \url{https://github.com/rafaellucian0/fraud_detection-big_data}, que inclui Docker Compose, configuracoes, jobs Spark, DAGs Airflow, testes e instrucoes para posicionar o dataset ULB.

\section{References}
\begin{thebibliography}{99}

\bibitem{akidau2015}
Akidau, T.; Chernyak, S.; Lax, R.; et al. (2015). The Dataflow Model: A Practical Approach to Balancing Correctness, Latency, and Cost in Massive-Scale, Unbounded, Out-of-Order Data Processing. \textit{Proceedings of the VLDB Endowment}, 8(12), 1792--1803.

\bibitem{adlermann2024}
Adlermann, J. F. (2024). A GRA-Enhanced Cloud AI Framework for Petabyte-Scale Multi-Tenant Environments: Multivariate Classification for Credit Card Fraud Detection and Adaptive Risk Analytics. \textit{International Journal of Engineering \& Extended Technologies Research}, 6(6).

\bibitem{allam2024}
Allam, K. (2024). Cloud-Based Data Hubs and SQL Pipelines for Real-Time Financial Analytics. \textit{International Journal of Emerging Trends in Computer Science and Information Technology}, 5(4), 94--104.

\bibitem{armbrust2021}
Armbrust, M.; Ghodsi, A.; Xin, R.; Zaharia, M. (2021). Lakehouse: A New Generation of Open Platforms that Unify Data Warehousing and Advanced Analytics. \textit{Proceedings of CIDR 2021}.

\bibitem{battapothu2025}
Battapothu, V. (2025). Lakehouse-Powered Payment Risk at Scale. \textit{Manuscrito tecnico independente}.

\bibitem{bfsilakehouse2023}
Oliveira, M.; Iyer, R. (2023). Migrating BFSI Data Workloads to Cloud-Native Environments: A Case Study on Multi-Tier Data Lakehouse Architectures with AWS Redshift, Athena, and Intelligent Orchestration for Compliance. \textit{Journal of Engineering, Mechanics and Modern Architecture}, 2(11), 50--61.

\bibitem{carrington2025}
Carrington, N. L. (2025). AI-Enabled Cloud Lakehouse for Large-Scale Data Warehousing: SAP Integration for Secure Analytics and Tableau-Driven Reporting. \textit{International Journal of Research and Applied Innovations}, 8(5).

\bibitem{dalpozzolo2015}
Dal Pozzolo, A.; Caelen, O.; Le Borgne, Y.-A.; Waterschoot, S.; Bontempi, G. (2015). Calibrating Probability with Undersampling for Unbalanced Classification. \textit{IEEE Symposium Series on Computational Intelligence}, 159--166.

\bibitem{dalpozzolo2018}
Dal Pozzolo, A.; Caelen, O.; Johnson, R. A.; Bontempi, G. (2018). Credit Card Fraud Detection: A Realistic Modeling and a Novel Learning Strategy. \textit{IEEE Transactions on Neural Networks and Learning Systems}, 29(8), 3784--3797.

\bibitem{davis2006}
Davis, J.; Goadrich, M. (2006). The Relationship Between Precision-Recall and ROC Curves. \textit{Proceedings of the 23rd International Conference on Machine Learning}, 233--240.

\bibitem{hand2009}
Hand, D. J. (2009). Measuring Classifier Performance: A Coherent Alternative to the Area Under the ROC Curve. \textit{Machine Learning}, 77, 103--123.

\bibitem{kreps2011}
Kreps, J.; Narkhede, N.; Rao, J. (2011). Kafka: A Distributed Messaging System for Log Processing. \textit{Proceedings of NetDB}, 1--7.

\bibitem{ngai2011}
Ngai, E. W. T.; Hu, Y.; Wong, Y. H.; Chen, Y.; Sun, X. (2011). The Application of Data Mining Techniques in Financial Fraud Detection: A Classification Framework and an Academic Review of Literature. \textit{Decision Support Systems}, 50(3), 559--569.

\bibitem{saito2015}
Saito, T.; Rehmsmeier, M. (2015). The Precision-Recall Plot Is More Informative than the ROC Plot When Evaluating Binary Classifiers on Imbalanced Datasets. \textit{PLOS ONE}, 10(3), e0118432.

\bibitem{ulb2013}
Worldline and Machine Learning Group, Universite Libre de Bruxelles. (2013). Credit Card Fraud Detection Dataset. \textit{Kaggle Dataset: mlg-ulb/creditcardfraud}.

\bibitem{zaharia2016}
Zaharia, M.; Xin, R. S.; Wendell, P.; et al. (2016). Apache Spark: A Unified Engine for Big Data Processing. \textit{Communications of the ACM}, 59(11), 56--65.

\end{thebibliography}
\end{document}
